\section{Exact Factor Anatomy and Pole Localization}
\subsection{Schur Closure}

The E05B holdout applied \eqref{eq:factor} to the 15 retained coalitions at 16
independent operating points. All builds were smooth-valid. The largest
pointwise log-determinant factorization residual was
$7.38964\times10^{-13}$; the largest residual after frequency integration was
$2.27374\times10^{-13}$; and the largest residual after finite M\"obius
differencing was $3.41061\times10^{-13}$. These values are consistent with
floating-point roundoff relative to the observed externalities.

The factorization answers a question that the aggregate characteristic
functional cannot: is a finite externality produced predominantly by
algebraic re-equilibration, action-dependent dynamic scaling, or pole motion?
Figure~\ref{fig:factor} reports the median holdout anatomy for four registered
examples.

\begin{figure}[t]
\centering
\includegraphics[width=\columnwidth]{figure_E05B_1_factor_anatomy.pdf}
\caption{Exact median holdout anatomy for registered coalitions. A7--A8 is
dominated by the algebraic factor, while A3--A6 has zero connected
algebraic/mass contribution and is pole dominated. Small triple coefficients
are retained rather than visually inflated.}
\label{fig:factor}
\end{figure}

For A7--A8, the median full externality was $-0.00411289$. Its algebraic,
mass, and pole components were $-0.00353228$, approximately zero, and
$-0.00059089$, respectively. The leading spectral interaction is therefore
largely an ac re-equilibration effect in the descriptor determinant, although
it also contains a reproducible pole component.

A3--A6 provides the complementary case. These two co-located action families
act on the bus-38 PLL and current-command bandwidths. Their connected
algebraic and mass terms are zero at the recorded precision, while the median
pole externality is $-0.00033610$. This is a direct counterexample to the idea
that every resolved log-determinant externality is merely an algebraic
power-flow artifact. Ten pairs and five triples exhibit reproducible
finite-pole externalities on the E05B holdout.

The triple results require careful language. A nonzero third-order coefficient
in $\Phi_{\mathrm{pole}}$ is a genuine third-order \emph{finite-pole
externality}: after removing all lower-order set contributions, the pole
factor remains non-additive. It does not show that three physical devices
form an isolated causal loop. Such a statement would require a frozen
localization object, a mechanism-specific intervention, and a matched negative
control, none of which is present here.

\subsection{Connected Logarithmic Derivative and Contour Moments}

Differentiating only the pole factor with respect to $s$ removes the
frequency-independent algebraic and mass factors:
\begin{equation}
\frac{\partial}{\partial s}\log\det(sI-A)
=\tr[(sI-A)^{-1}].
\label{eq:logderiv}
\end{equation}
Applying the finite connected difference to \eqref{eq:logderiv} gives a
resolvent-based description of non-additive pole displacement. Registered
closed contours then provide zeroth- and first-order moments. A certificate is
accepted only when the contour stays clear of poles at every required action
vertex, occupancy is consistent, quadrature refinement closes, and direct and
connected moments agree.

Three candidates met these conditions across all 16 holdout points: the pair
A7--A8 and triples A2--A7--A8 and A2--A5--A8. For A7--A8, the median absolute
first-moment shift was $4.74913\times10^{-3}\ \mathrm{s}^{-1}$; the median
real and imaginary components were $4.65632\times10^{-3}$ and
$-8.99434\times10^{-4}\ \mathrm{s}^{-1}$, respectively. Maximum numerical
contour error was $1.24\times10^{-15}$. Both triple candidates likewise
retained valid one-pole certificates and same-sign fraction 1.0.

The A3--A6 contour did not pass. Its maximum numerical contour error was
0.9151, leaving zero valid-resolved holdout points under the registered rule.
This negative case is retained. The factor decomposition still proves that
A3--A6 has a pole externality; it does not authorize a contour-localized pole
claim for that coalition.

\subsection{Relation to Return-Ratio, Impedance, and Spectral Methods}

Return-difference matrices and impedance Nyquist tests diagnose feedback
stability and interconnection effects
\cite{macfarlane1970return,sun2011impedance,rygg2016sequence}. They can often
be measured or assembled with fewer internal model details than
\eqref{eq:descriptor}. Proximal eigenvalue-displacement methods identify modes
whose coupling invalidates a first-order open-loop approximation
\cite{dewangan2026interaction}. That recent comparator returns converter
interaction modes from open-loop proximity and first-/second-order eigenvalue
displacement. It overlaps with the present interest in modal interaction, but
does not define a finite action-coalition M\"obius estimand, require a new AC
equilibrium at every subset vertex, separate descriptor-factor externalities,
or test the independently identified minimum-damping family. The present
method is neither its replacement nor a new stability criterion.

Its output is instead an attribution on a finite action set. The same physical
portfolio is evaluated at every subset, and inclusion--exclusion isolates what
belongs only to the joint finite action. The connected formula then states how
that set-function coefficient is assembled from total operator derivatives.
Finally, \eqref{eq:factor} asks whether the coefficient lies in algebraic
re-equilibration, mass scaling, or finite poles. These questions are
complementary to loop stability and modal proximity.

The determinant language also requires restraint. In finite non-normal
matrices, log determinants and contour moments are computational identities
under their stated regularity conditions. They do not inherit the ordering,
positivity, or causal interpretation of self-adjoint spectral-shift theory.
This study therefore uses \emph{connected pole localization}, not spectral
causality.
